% !TEX root = ../main.tex
%%
% BIThesis 本科毕业设计论文模板 —— 使用 XeLaTeX 编译
%%

% 第二章节
\chapter{相关技术与方法}
\section{软件测试相关技术}
\subsection{单元测试与集成测试}
单元测试（Unit Testing）是软件测试体系中最基础的测试级别，其核心目标是对软件中的最小可测试单元进行验证~\cite{Myers2011}。在面向对象编程范式下，单元通常指一个函数、一个类或一组紧密相关的函数集合。单元测试具有以下典型特征：每个测试用例应具有原子性，专注于验证单一功能点；测试用例之间应相互独立，避免共享状态导致测试结果的相互影响；单元测试应在隔离环境中执行，通过模拟（Mock）或桩（Stub）技术排除对外部依赖的依赖。单元测试的优势在于能够快速定位缺陷，降低调试成本，同时也是代码重构的重要保障。

集成测试（Integration Testing）旨在验证多个软件模块或组件之间交互的正确性~\cite{Myers2011}。与单元测试不同，集成测试关注的是模块接口、调用协议、数据交换及集成后的系统行为。集成测试能够发现模块接口不匹配、数据格式不一致、调用时序错误等问题，是单元测试与系统测试之间的重要桥梁。在持续集成（Continuous Integration, CI）环境中，集成测试通常配置为自动化执行，及时发现代码集成过程中引入的缺陷。

在回归测试场景中，单元测试与集成测试具有互补作用。单元测试用于验证代码变更对单个函数或类的影响，侧重于局部正确性验证；集成测试则用于验证代码变更对模块间协作的影响，侧重于全局行为验证~\cite{YooHarman2012}。本研究聚焦于面向代码变更的回归测试用例生成，需要同时考虑单元测试级别与集成测试级别的用例生成。

\subsection{自动化测试框架（Pytest）}
Pytest是Python生态中应用最为广泛的自动化测试框架之一，以其简洁的API设计、强大的插件生态及丰富的功能特性著称~\cite{pytestOfficial}。相较于Python标准库中的unittest框架，Pytest采用更加灵活的设计理念，支持基于函数的测试用例定义，降低了测试代码的编写门槛。Pytest的核心设计原则包括：约定优于配置、失败即停、以及详细的断言重写。

Pytest的Fixture机制是其最具特色的功能之一，Fixture用于提供测试所需的上下文环境或依赖资源，支持多种作用域（function、class、module、session）及参数化配置~\cite{pytestOfficial}。此外，社区提供pytest-asyncio、pytest-mock、pytest-xdist等大量插件。本研究基于Pytest构建测试执行环境。

\section{代码变更分析技术}
\subsection{Unified Diff解析}
代码变更分析的第一步是对版本控制系统产生的差异文件进行解析。Git作为目前最流行的分布式版本控制系统，其差异输出格式（Diff Format）是代码变更分析领域的事实标准~\cite{ChaconStraub2014}。Git Diff采用Unified Diff格式进行展示，新增行以前缀\textsf{+}标识，删除行以前缀\textsf{-}标识，上下文行以空格前缀标识。

Python标准库中的\textsf{difflib}模块提供了差异计算的基础功能~\cite{PythonDifflib}。为实现符合Unified Diff规范的解析与处理，社区广泛使用Python包\textsf{unidiff}，该库提供了从统一差异格式到结构化数据模型的双向转换能力~\cite{unidiffLibrary}。通过\textsf{unidiff}，开发者可将原始Diff文本解析为包含\textsf{PatchedFile}、\textsf{SourceFile}、\textsf{Hunk}等对象的数据结构，每个对象封装了文件路径、变更块起止位置、具体变更行等属性信息。

Unified Diff解析在实际应用中面临若干技术挑战：变更块的重叠与合并问题、二进制文件与文件重命名的处理、编码与换行符问题等。本研究在代码变更解析模块中集成了Unified Diff解析功能，用于将Git提交产生的差异文件转换为结构化的变更表示。

\subsection{代码变更图与依赖关系建模}
代码变更的影响分析需要借助程序结构模型来追踪变更的传播路径。程序依赖图（Program Dependence Graph, PDG）是这一领域最具代表性的表示模型~\cite{Ferrante1987}。PDG将程序的控制流依赖与数据流依赖统一建模，通过PDG可以精确追踪程序中任意语句对其他语句的影响范围。

调用图（Call Graph）是描述函数或方法调用关系的重要模型~\cite{Grove1995CallGraph}。在调用图中，节点表示函数实体，有向边表示调用关系。调用图可分为静态调用图和动态调用图两种类型。静态调用图通过解析代码中的函数调用表达式构建，覆盖所有可能的调用路径；动态调用图则记录程序实际执行过程中的调用序列，覆盖范围有限但准确性更高。在变更影响分析中，调用图支持自底向上的影响传播。

除调用关系外，依赖关系还包括模块依赖、类依赖及数据依赖等多种类型~\cite{AnselmGoldberg2003}。本研究构建了综合依赖模型，整合调用图、控制流依赖与数据流依赖等多种关系表示，通过图遍历算法实现变更影响范围的有效识别。

影响分析方法可分为基于覆盖率的测试选择与基于传播的变更影响分析两类~\cite{BohnerArnold1996}。基于覆盖率的方法通过比较变更前后的代码覆盖率差异来确定受影响的测试用例集~\cite{Chilakapati1998}。基于传播的方法则从变更点出发，沿程序依赖关系图进行传播，识别所有直接或间接受影响的程序实体~\cite{Le2018ImpactPropagation}。本研究综合采用两类方法。

\subsection{代码语义表示方法（AST与LLM）}
代码的结构化表示是进行语义级代码分析的基础。抽象语法树（Abstract Syntax Tree, AST）将源代码转换为树形结构，其中每个节点对应源代码中的一个语法构造~\cite{Aho2006Compilers}。AST在保留代码语法结构的同时移除了无关的格式信息，使得程序分析可以在结构化层面进行。Python的AST模块可将源代码解析为\textsf{ast.Module}根节点下的层次结构~\cite{PythonAST}。

Tree-sitter是一个用于构建多语言解析器的增量式解析系统~\cite{Brunsfeld2018TreeSitter}。与传统的解析器不同，Tree-sitter具有语言无关的解析框架、增量解析能力、错误容忍机制等特点，使其特别适用于代码变更分析场景。

大语言模型（Large Language Models, LLMs）的出现为代码语义表示提供了新的技术路径~\cite{Fan2023LLMSurvey}。预训练代码模型（如Codex、CodeT5、StarCoder等）在大规模代码语料上进行预训练，学习到了代码的语义表示与生成能力~\cite{Chen2021Codex}。代码嵌入（Code Embedding）将代码片段映射为稠密向量表示，使得代码之间的语义相似性可通过向量距离进行度量~\cite{Hindle2012Naturalness}。

然而，纯文本表示的LLM在处理结构化程序代码时存在固有限制。常见思路包括基于AST的序列化、基于图神经网络的程序图编码以及多模态信息融合等。本研究结合Tree-sitter的结构化解析与LLM的语义理解，实现从语法节点到语义推理的衔接。

\section{多智能体系统技术}
\subsection{基于LangGraph的工作流建模}
LangGraph是LangChain生态中用于构建有状态、多角色交互应用的工作流编排框架~\cite{LangGraphOfficial}。与传统的线性链条式（Chain）处理流程不同，LangGraph采用图结构对复杂工作流进行建模，其中节点（Node）表示处理单元，边（Edge）表示控制流与数据流。这种图结构天然支持条件分支、循环迭代及多路径并行等复杂控制模式。

LangGraph的核心抽象包括StateGraph、节点（Node）和边（Edge）三个层次：状态模式常借助Pydantic等机制描述全局状态~\cite{PydanticDoc}；节点以函数形式实现状态的读入与更新；边刻画普通迁移、条件分支及入口/终止等控制关系。框架支持惰性求值与状态持久化。

本研究采用LangGraph作为多智能体协同框架的核心基础设施。通过StateGraph定义测试生成工作流的全局状态结构，工作流中的各个处理阶段均建模为独立节点，通过边定义节点间的数据流向与控制依赖。这种基于图的工作流建模方式既保证了处理流程的清晰性，又为后续的功能扩展与优化提供了灵活性。

\subsection{多Agent协同机制}
多智能体系统（Multi-Agent System, MAS）是由多个具有自主能力的智能体（Agent）组成的计算系统，智能体之间通过协作完成复杂任务~\cite{Wooldridge2009MAS}。与单一智能体相比，多智能体系统在问题分解、信息集成、容错冗余及并行处理等方面具有显著优势~\cite{Stone2000MultiAgent}。

智能体的角色定义是多智能体系统设计的基础~\cite{Kraus2008AgentRoles}。每个智能体通常被赋予特定的角色，承担特定的功能职责。在软件工程应用场景中，常见的角色包括分析型Agent（负责理解需求、分析问题）、规划型Agent（负责制定执行计划、分解任务）、执行型Agent（负责具体操作的实施）、验证型Agent（负责检查结果正确性）及协调型Agent（负责管理多个Agent的协作流程）~\cite{Wu2023AutoGen}。

信息共享是多智能体协同的核心机制~\cite{Franklin1996AgentComm}，常见模式包括点对点消息、黑板、发布-订阅与消息队列等。在基于LLM的系统中，常通过共享上下文或对话历史传递信息~\cite{Hong2023MetaGPT}。对于可并行的子任务可并行执行，存在依赖的子任务则按拓扑顺序串行~\cite{Fagin2003StateSync}。当多个Agent结论不一致时，可采用权威裁定、投票、置信度比较或协商等策略~\cite{Ossowski2014Consensus}。

本研究构建的多智能体测试生成系统包含11个专业Agent，分别负责代码变更解析、AST分析、影响范围推理、测试策略规划、测试用例生成、测试执行、结果评估、缺陷模式匹配、代码检索、上下文管理和报告生成等功能。

\section{大语言模型技术}
\subsection{代码理解与结构化生成}
大语言模型在代码理解任务中的能力源于其在大规模代码语料上的预训练~\cite{Fan2023LLMSurvey}。预训练阶段，模型学习到编程语言的语法规则、语义模式及常见的编码范式~\cite{Feng2020CodeBERT}。代码理解任务包括代码摘要、代码搜索、代码补全及代码翻译等~\cite{Wang2021CodeT5}。

代码嵌入将代码映射为稠密向量，使语义相近片段在向量空间中彼此接近~\cite{Alon2019Code2Vec}。

结构化生成要求模型输出符合既定代码形态与项目规范。常见对策包括约束解码、提示工程与生成后校验等~\cite{Geng2023ConstrainedDecoding,Bommasani2021FoundationModels}。本研究结合提示模板与生成后检查，以约束测试代码形态。

检索增强生成（RAG）通过从外部库检索相关片段并拼入提示，扩展模型可用知识~\cite{Lewis2020RAG}。本研究在流水线中集成了这一机制。

\subsection{测试用例生成与修复}
基于大语言模型的测试用例生成已成为软件测试领域的研究热点~\cite{Tufano2021UnitTest}。其核心难点包括：如何把变更信息组织为有效提示、如何约束输出符合工程规范、如何评价充分性等。

提示设计需要兼顾任务说明、上下文、格式约束与示例等要素~\cite{Brown2020GPT3}；也可采用思维链/分步规划等提示策略以提升筹划性。

测试用例修复用于应对生成中的语法与语义问题~\cite{Xia2023TestRepair}，既可通过静态分析修补，也可结合执行反馈迭代。Self-healing等研究总结了失败模式并给出自动修复策略~\cite{Yin2021SelfHealingTest}。

充分性常用覆盖率等指标刻画；亦可通过故障注入考察检出能力。本研究以覆盖率与执行成功率为主，并在失败时进入修复回路。

综上所述，本章介绍了与本研究相关的核心技术与方法。软件测试技术为研究提供了测试理论与工程实践基础；代码变更分析技术为影响范围推理提供了技术支撑；多智能体系统技术为复杂测试生成任务的协同处理提供了架构框架；大语言模型技术为代码理解与测试用例生成提供了智能能力。这些技术与方法相互结合，共同构成本研究的技术基础。
